% =============================================================================
% Worked solutions for Listas/PS1.pdf (Corporate Finance -- Problem Set 1).
% Short answer key: Listas/PS1-solutions-short.tex.
% Compile: pdflatex (twice); check: pdffonts (Type 1 only, uni = yes).
% Numbers: python -I Listas/PS1-solutions_check.py  -> PASS
% =============================================================================
\documentclass[11pt,a4paper]{article}
\usepackage[utf8]{inputenc}
\usepackage[T1]{fontenc}
\usepackage{lmodern}   % Type 1 vector fonts (replaces PK bitmaps)
\usepackage{cmap}      % ToUnicode CMap -> copyable text and accents
\usepackage[english]{babel}
\usepackage{amsmath,amssymb,amsfonts,amsthm}
\usepackage[a4paper,margin=2.2cm]{geometry}
\usepackage{booktabs}
\usepackage{enumitem}
\usepackage{fancyhdr}
\usepackage{xcolor}
\usepackage{microtype}

\pagestyle{fancy}\fancyhf{}
\lhead{\footnotesize Finan\c{c}as II -- Corporativas (MPE Insper)}
\rhead{\footnotesize Problem Set 1 --- Worked solutions}
\cfoot{\footnotesize\thepage}
\renewcommand{\headrulewidth}{0.3pt}
\setlength{\headheight}{14pt}

\makeatletter
\renewcommand{\@seccntformat}[1]{\csname the#1\endcsname.\quad}
\makeatother
\numberwithin{equation}{section}
\setlist{topsep=2pt,itemsep=1pt}

\newcommand{\step}[2]{\par\medskip\noindent\textbf{Step #1 --- #2.}~}
\newenvironment{notation}{\par\noindent\textbf{Notation.}\begin{itemize}[leftmargin=1.4em]}{\end{itemize}}
\newcommand{\intuition}[1]{\par\medskip\noindent\textbf{Intuition.}~#1}
\newcommand{\refline}[1]{\par\smallskip\noindent{\small\color{black!70}\textbf{Ref:} #1}}
\newcommand{\remark}[1]{\par\smallskip\noindent\textit{Remark.}~#1}
\newcommand{\answer}[1]{\par\medskip\noindent\fbox{\parbox{\dimexpr\linewidth-2\fboxsep-2\fboxrule}{\textbf{Answer.}~#1}}\par}

\begin{document}

\begin{center}
  {\Large\bfseries Corporate Finance --- Problem Set 1}\\[4pt]
  {\large Worked solutions}
\end{center}

\noindent Statement: \texttt{Listas/PS1.pdf}. Each exercise is restated briefly, then solved
step by step. Rates are written as decimals inside equations. References are to Berk, DeMarzo \&
Harford, \emph{Fundamentals of Corporate Finance}, 2nd ed.\ (2011), printed pages. Every number
below is reproduced by \texttt{Listas/PS1-solutions\_check.py}. A one-page answer key is in
\texttt{Listas/PS1-solutions-short.pdf}.

% =============================================================================
\section{Exercise 1 --- Net income, free cash flow and deferred taxes}

\begin{notation}
  \item $NI_t$: net income of year $t$; $EBIT_t$: earnings before interest and taxes.
  \item $EBITDA_t = EBIT_t + Dep_t$: operating profit before depreciation and amortization.
  \item $Dep_t\ge 0$: depreciation (a non-cash expense); $CAPEX_t\ge0$: capital expenditures.
  \item $NWC_t$: net working capital (operating current assets minus non-interest-bearing current
        liabilities); $\Delta NWC_t = NWC_t - NWC_{t-1}$, any real number.
  \item $Int_t\ge0$: interest expense; $\tau_c\in[0,1)$: corporate tax rate.
  \item $FCF_t$: unlevered free cash flow of year $t$, cash generated by the assets for
        \emph{all} investors.
  \item $T^{book}_t$: income tax expense on the income statement; $T^{cash}_t$: taxes actually paid.
  \item $DTL_t\ge0$: deferred tax liability on the balance sheet at the end of year $t$;
        $\Delta DTL_t = DTL_t - DTL_{t-1}$.
\end{notation}

\subsection*{(a) Why discount FCF and not net income?}

\step{1}{What a valuation must measure} The value of an asset is the present value of the
\emph{cash} it delivers to its investors. Investors can only consume or reinvest cash, and the time
value of money depends on \emph{when} the cash arrives. A discount factor $1/(1+r)^t$ is a price of
a dollar received at date $t$, so the object multiplied by it must be dollars received at date $t$.
Net income is an accrual measure: it allocates revenues and costs to periods following accounting
conventions (matching principle), not to the dates on which cash moves.

\step{2}{The textbook FCF} Free cash flow is built from operating profit, as if the firm had no
debt:
\begin{equation}\label{eq:1fcf}
  FCF_t = EBIT_t(1-\tau_c) + Dep_t - CAPEX_t - \Delta NWC_t .
\end{equation}

\step{3}{Net income in the same terms} Net income deducts interest before taxes:
\begin{equation}\label{eq:1ni}
  NI_t = (EBIT_t - Int_t)(1-\tau_c).
\end{equation}
Solving \eqref{eq:1ni} for $EBIT_t(1-\tau_c)$:
\begin{equation}\label{eq:1ebit}
  EBIT_t(1-\tau_c) = NI_t + (1-\tau_c)\,Int_t .
\end{equation}

\step{4}{The bridge from NI to FCF} Substituting \eqref{eq:1ebit} into \eqref{eq:1fcf}:
\begin{equation}\label{eq:1bridge}
  FCF_t = \bigl(NI_t + (1-\tau_c)\,Int_t\bigr) + Dep_t - CAPEX_t - \Delta NWC_t .
\end{equation}

\step{5}{The four differences} Reading \eqref{eq:1bridge} term by term:
\begin{enumerate}[label=(\roman*),leftmargin=2em]
  \item \textbf{Non-cash expenses} ($+Dep_t$). Depreciation and amortization reduce $NI$ but no
        cash leaves the firm. They matter for value only through the taxes they save,
        $\tau_c Dep_t$, which \eqref{eq:1fcf} keeps.
  \item \textbf{Capital expenditures} ($-CAPEX_t$). The purchase of equipment is cash out
        \emph{today}, but the income statement spreads it over the asset's life as depreciation.
        $NI$ therefore shows the cost too late; FCF shows it when it is paid.
  \item \textbf{Working capital} ($-\Delta NWC_t$). Revenue is booked when the sale is made, even if
        the customer pays later (receivables); inventory is paid for before the good is sold; part
        of costs is paid later (payables). An increase in $NWC$ is cash tied up in operations
        and is subtracted; a decrease releases cash.
  \item \textbf{Financing flows} ($+(1-\tau_c)Int_t$). $NI$ is the residual for shareholders,
        after interest. FCF is the cash available to \emph{all} capital providers, so it is
        computed before any payment to them. The financing side (interest, its tax shield) is
        handled in the discount rate (WACC) or valued separately (APV). Discounting $NI$ at the
        WACC would mix an equity cash flow with a firm-wide discount rate and count debt twice.
\end{enumerate}

\step{6}{A numerical illustration (data of Exercise 2)} In Exercise 2 below, the corrected project
earns $NI = 5{,}525$ per year for 10 years (thousands of dollars), total $55{,}250$; its FCF are
$-35{,}000$ at $t=0$, $8{,}025$ in years 1--9 and $18{,}025$ in year 10, total
\begin{equation}\label{eq:1sum}
  -35{,}000 + 10\times 8{,}025 + 10{,}000 = 55{,}250 .
\end{equation}
Over the whole life the two \emph{totals} coincide: depreciation adds up to capex and the working
capital comes back. Only the \emph{timing} differs. At 14\% the present value of the NI stream is
$5{,}525\times 5.216116 = 28{,}819.04$, while the NPV of the FCF stream is $9{,}556.77$ (computed
in Section~2). Discounting NI would overvalue the project by about $19{,}262$ because it pushes the
25 million of capex into the future and ignores the 10 million of working capital tied up for ten
years.

\answer{We discount FCF because value is the present value of cash received by investors, dated
when it is received. Net income differs from FCF because it (i) deducts non-cash depreciation,
(ii) ignores capital expenditures when they are paid, (iii) ignores investment in working capital
(accrual revenues and costs), and (iv) is net of interest, i.e.\ it mixes financing with operating
flows. Formally $FCF = NI + (1-\tau_c)Int + Dep - CAPEX - \Delta NWC$.}

\refline{Berk, \S2.6 (The statement of cash flows), pp.~42--44; \S9.3 (Determining incremental
free cash flow: converting from earnings to free cash flow), pp.~254--259.}

\subsection*{(b) Should FCF be adjusted for changes in deferred taxes?}

\step{1}{Book taxes versus cash taxes} The statement gives the accounting identity
\begin{equation}\label{eq:1tax}
  T^{book}_t = T^{cash}_t + \Delta DTL_t .
\end{equation}
Solving \eqref{eq:1tax} for the cash actually paid:
\begin{equation}\label{eq:1cash}
  T^{cash}_t = T^{book}_t - \Delta DTL_t .
\end{equation}

\step{2}{Which tax enters the textbook formula} In \eqref{eq:1fcf}, $EBIT_t(1-\tau_c) = EBIT_t -
\tau_c EBIT_t$: the tax deducted, $\tau_c EBIT_t$, is the \emph{book} tax on operating income,
computed from book EBIT (book depreciation). Write it as $T^{book,op}_t = \tau_c EBIT_t$.

\step{3}{FCF with cash taxes} The cash flow the assets really release subtracts the taxes really
paid on operating income, $T^{cash,op}_t$:
\begin{equation}\label{eq:1fcfcash}
  FCF^{cash}_t = EBIT_t - T^{cash,op}_t + Dep_t - CAPEX_t - \Delta NWC_t .
\end{equation}
Applying \eqref{eq:1cash} to the operating taxes, $T^{cash,op}_t = \tau_c EBIT_t - \Delta DTL^{op}_t$.
Substituting into \eqref{eq:1fcfcash}:
\begin{equation}
  FCF^{cash}_t = EBIT_t - \bigl(\tau_c EBIT_t - \Delta DTL^{op}_t\bigr) + Dep_t - CAPEX_t - \Delta NWC_t , \notag
\end{equation}
and collecting terms,
\begin{equation}\label{eq:1adj}
  FCF^{cash}_t = EBIT_t(1-\tau_c) + Dep_t + \Delta DTL^{op}_t - CAPEX_t - \Delta NWC_t .
\end{equation}
So yes: the increase in the (operating) deferred tax liability is added back, exactly like
depreciation. It is an expense on the income statement that is not paid in cash this period. This
is how the statement of cash flows treats it as well (Berk, p.~42: ``we add back any other non-cash
expenses (for example, deferred taxes)'').

\step{4}{A two-year example} An asset costing 100 yields $EBITDA=100$ in each of two years,
$\tau_c=0.30$, $r=0.10$. Books depreciate straight line (50, 50); the tax code lets the firm
expense it all in year 1 (100, 0). Then book $EBIT=50$ in both years, book tax $0.30\times50=15$,
and the textbook operating cash flow from \eqref{eq:1fcf} (capex aside) is
\begin{equation}\label{eq:1ex85}
  50(1-0.30) + 50 = 85 \quad\text{in each year.}
\end{equation}
Cash taxes are $0.30\times(100-100)=0$ in year 1 and $0.30\times(100-0)=30$ in year 2. By
\eqref{eq:1cash}, $\Delta DTL_1 = 15-0 = +15$ and $\Delta DTL_2 = 15-30 = -15$. With \eqref{eq:1adj}:
\begin{equation}\label{eq:1ex100}
  FCF^{cash}_1 = 85 + 15 = 100, \qquad FCF^{cash}_2 = 85 - 15 = 70,
\end{equation}
which equal $EBITDA - T^{cash}$ directly ($100-0$ and $100-30$). The present values are
\begin{align}
  PV^{book} &= \frac{85}{1.10} + \frac{85}{(1.10)^2} = 77.2727 + 70.2479 = 147.5207, \notag\\
  PV^{cash} &= \frac{100}{1.10} + \frac{70}{(1.10)^2} = 90.9091 + 57.8512 = 148.7603. \label{eq:1expv}
\end{align}
The difference, $148.7603-147.5207=1.2397 = 15/1.10 - 15/1.10^2$, is the value of paying 15 of
taxes one year later. Ignoring $\Delta DTL$ misses it.

\step{5}{How to apply it in practice} Four caveats:
\begin{enumerate}[label=(\roman*),leftmargin=2em]
  \item \textbf{Only operating deferred taxes.} Adjust for deferred taxes created by operating items
        (typically accelerated tax depreciation versus straight-line book depreciation). Deferred
        taxes from non-operating items (e.g.\ pension or financial gains) do not belong in operating FCF.
        A deferred tax \emph{asset} enters with the opposite sign: its increase is subtracted.
  \item \textbf{Consistency with the equity bridge.} If $\Delta DTL$ is added to FCF, the stock of
        DTL must \emph{not} also be subtracted from enterprise value as a debt-like claim. Doing both
        counts the same deferral twice. The alternative treatment (DTL as quasi-debt, no
        $\Delta DTL$ in FCF) is also consistent; mixing the two is not.
  \item \textbf{Reversal and the terminal value.} For a single asset the timing difference reverses:
        in the example $\Delta DTL_1 + \Delta DTL_2 = 0$. Only the time value is gained. For a
        growing firm that keeps investing, new deferrals offset reversals and DTL grows with the
        asset base; the deferral becomes close to permanent. In the terminal year, project
        $\Delta DTL$ in line with growth, not as a one-off spike.
  \item \textbf{Equivalent route in project analysis.} In capital budgeting (as in Berk, Ch.~9),
        taxes are computed with the \emph{tax} depreciation schedule (e.g.\ MACRS). The formula
        $FCF = (EBITDA - Dep^{tax})(1-\tau_c) + Dep^{tax} - CAPEX - \Delta NWC$ then gives cash taxes
        directly and no separate $\Delta DTL$ term is needed. In the example it gives
        $(100-100)(0.70)+100 = 100$ and $(100-0)(0.70)+0 = 70$, the same as \eqref{eq:1ex100}.
\end{enumerate}

\answer{Yes. Valuation must use taxes actually paid. Since $T^{cash}=T^{book}-\Delta DTL$, the
textbook formula should add the increase in operating deferred tax liabilities:
$FCF_t = EBIT_t(1-\tau_c) + Dep_t + \Delta DTL_t - CAPEX_t - \Delta NWC_t$. Deferred taxes are a
non-cash expense, like depreciation. Keep it to operating deferred taxes, do not also treat the DTL
as debt in the equity bridge, and model its reversal or steady growth in the terminal value.
Equivalently, compute taxes on EBIT with tax depreciation.}

\refline{Berk, \S2.6, p.~42 (deferred taxes added back as a non-cash expense); \S9.3, pp.~254--259;
\S9.4 (Adjusting free cash flow: accelerated depreciation), pp.~262--264.}

% =============================================================================
\section{Exercise 2 --- Free cash flows and NPV of a new equipment}

\noindent\textbf{Data (thousands of dollars).} Equipment $25{,}000$ bought at $t=0$; consultancy
fee $1{,}000$ already paid; the report shows, for each year $t=1,\dots,10$: sales $30{,}000$,
COGS $18{,}000$, overhead $2{,}000$, depreciation $2{,}500$, operating income $7{,}500$, income tax
$2{,}625$, net income $4{,}875$. Working capital $10{,}000$ at $t=0$, fully recovered at $t=10$. Only
$1{,}000$ of the $2{,}000$ overhead is incremental. Cost of capital 14\%.

\begin{notation}
  \item $I_0 = 25{,}000$: capex at $t=0$; $Dep = 2{,}500$: annual straight-line depreciation.
  \item $S=30{,}000$: sales; $C=18{,}000$: COGS; $O^{inc}$: incremental overhead per year.
  \item $W = 10{,}000$: working capital invested at $t=0$ and recovered at $t=10$.
  \item $\tau_c$: tax rate (to be inferred); $r=0.14$: cost of capital.
  \item $FCF_t$, $t=0,\dots,10$: incremental free cash flow; $\mathrm{NPV}$: value of the project at $t=0$.
  \item $A(r,T) = \bigl(1-(1+r)^{-T}\bigr)/r$: present value of 1 per year for $T$ years.
\end{notation}

\subsection*{(a) Free cash flows, years 0 to 10}

\step{1}{Infer the tax rate} The report taxes operating income of $7{,}500$ at $2{,}625$:
\begin{equation}\label{eq:2tau}
  \tau_c = \frac{2{,}625}{7{,}500} = 0.35 .
\end{equation}
Depreciation is consistent with the data: $Dep = 25{,}000/10 = 2{,}500$, so the equipment is fully
depreciated after 10 years; with no salvage value given, it is worth zero at $t=10$.

\step{2}{What is \emph{not} incremental} Only cash flows that change because the project is
accepted enter the analysis.
\begin{itemize}[leftmargin=1.4em]
  \item The \textbf{consultancy fee} (1 million) is a \emph{sunk cost}: it has been paid whatever we
        decide. It is excluded.
  \item \textbf{Overhead}: 1 million of the 2 million attributed is incurred even if the project is
        rejected. It is an allocation, not an incremental cost. Only $O^{inc}=1{,}000$ enters.
\end{itemize}

\step{3}{Corrected incremental EBIT} With $O^{inc}=1{,}000$:
\begin{align}
  EBIT &= S - C - O^{inc} - Dep \notag\\
       &= (30{,}000) - (18{,}000) - (1{,}000) - (2{,}500) \notag\\
       &= 8{,}500 . \label{eq:2ebit}
\end{align}

\step{4}{Unlevered net income (NOPAT)} Using \eqref{eq:2tau} and \eqref{eq:2ebit}:
\begin{align}
  EBIT(1-\tau_c) &= (8{,}500)(1-0.35) \notag\\
                 &= 8{,}500 - 2{,}975 \notag\\
                 &= 5{,}525 . \label{eq:2nopat}
\end{align}

\step{5}{FCF in years 1--9} No capex and no change in working capital in these years, so by
\eqref{eq:1fcf} and \eqref{eq:2nopat}:
\begin{align}
  FCF_t &= (5{,}525) + (2{,}500) - 0 - 0 \notag\\
        &= 8{,}025, \qquad t=1,\dots,9. \label{eq:2fcf19}
\end{align}

\step{6}{FCF at $t=0$} Capex and the working-capital build-up ($\Delta NWC_0 = +10{,}000$):
\begin{align}
  FCF_0 &= -(25{,}000) - (10{,}000) \notag\\
        &= -35{,}000 . \label{eq:2fcf0}
\end{align}

\step{7}{FCF at $t=10$} The operating flow \eqref{eq:2fcf19} plus the recovery of working capital
($\Delta NWC_{10} = -10{,}000$, i.e.\ $+10{,}000$ of cash):
\begin{align}
  FCF_{10} &= (8{,}025) - (-10{,}000) \notag\\
           &= 18{,}025 . \label{eq:2fcf10}
\end{align}

\begin{center}
\begin{tabular}{lrrr}
  \toprule
  (thousands of dollars)      & $t=0$       & $t=1,\dots,9$ & $t=10$ \\
  \midrule
  Sales                       &             & 30{,}000      & 30{,}000 \\
  $-$ COGS                    &             & $-18{,}000$   & $-18{,}000$ \\
  $-$ Incremental overhead    &             & $-1{,}000$    & $-1{,}000$ \\
  $-$ Depreciation            &             & $-2{,}500$    & $-2{,}500$ \\
  $=$ EBIT                    &             & 8{,}500       & 8{,}500 \\
  $-$ Taxes (35\%)            &             & $-2{,}975$    & $-2{,}975$ \\
  $=$ Unlevered net income    &             & 5{,}525       & 5{,}525 \\
  $+$ Depreciation            &             & 2{,}500       & 2{,}500 \\
  $-$ Capex                   & $-25{,}000$ &               &  \\
  $-\Delta$NWC                & $-10{,}000$ &               & $+10{,}000$ \\
  \midrule
  $=$ FCF                     & $-35{,}000$ & 8{,}025       & 18{,}025 \\
  \bottomrule
\end{tabular}
\end{center}

\answer{$FCF_0=-35{,}000$; $FCF_1=\dots=FCF_9=8{,}025$; $FCF_{10}=18{,}025$ (thousands of dollars).}

\remark{An equivalent shortcut is $FCF_t = (S-C-O^{inc})(1-\tau_c) + \tau_c Dep$ for operating
years: $(11{,}000)(0.65) + (0.35)(2{,}500) = 7{,}150 + 875 = 8{,}025$. The second term is the
depreciation tax shield.}

\subsection*{(b) Value of the project at 14\%}

\step{1}{Annuity factor} With $r=0.14$ and $T=10$:
\begin{align}
  (1.14)^{10} &= 3.707221, \notag\\
  A(0.14,10) &= \frac{1-(3.707221)^{-1}}{0.14} = \frac{1-0.269744}{0.14} = 5.216116 . \label{eq:2af}
\end{align}

\step{2}{NPV} Write $FCF_{10}$ as the regular $8{,}025$ plus the working capital $10{,}000$. Then
\begin{equation}\label{eq:2npvdef}
  \mathrm{NPV} = FCF_0 + 8{,}025\,A(0.14,10) + \frac{10{,}000}{(1.14)^{10}} .
\end{equation}
Substituting \eqref{eq:2fcf0} and \eqref{eq:2af} into \eqref{eq:2npvdef}:
\begin{align}
  \mathrm{NPV} &= (-35{,}000) + (8{,}025)(5.216116) + (10{,}000)(0.269744) \notag\\
               &= -35{,}000 + 41{,}859.33 + 2{,}697.44 \notag\\
               &= 9{,}556.77 . \label{eq:2npv}
\end{align}

\answer{The project is worth $\mathrm{NPV}\approx 9{,}556.8$ thousand, about \$9.56 million.
Accept it. (The present value of its future cash flows is $41{,}859.33+2{,}697.44=44{,}556.77$
thousand; the 35 million outlay is deducted from it.)}

\step{3}{Verification: from the consultants' net income to the NPV} The table below starts from
discounting the consultants' net income ($4{,}875$ per year) and adds, one at a time, each
correction, every row at 14\%:
\begin{center}
\begin{tabular}{lr}
  \toprule
  (thousands of dollars, PV at 14\%)                                         & Value \\
  \midrule
  PV of consultants' net income, $4{,}875\times5.216116$                     & 25{,}428.56 \\
  $+$ non-incremental overhead after tax, $1{,}000(1-0.35)\times5.216116$    & 3{,}390.48 \\
  $+$ depreciation add-back, $2{,}500\times5.216116$                         & 13{,}040.29 \\
  $-$ capex at $t=0$                                                         & $-25{,}000.00$ \\
  $-$ working capital, $10{,}000 - 10{,}000\times0.269744$                   & $-7{,}302.56$ \\
  \midrule
  $=$ NPV                                                                     & 9{,}556.77 \\
  \bottomrule
\end{tabular}
\end{center}
The sum reproduces \eqref{eq:2npv}. The sunk consultancy fee appears nowhere. The IRR of the FCF
stream is about 20.2\%, above 14\%, as it must be when the NPV is positive.

\intuition{Two corrections raise value (the overhead that would be paid anyway, and the
depreciation that is not cash); two lower it (the equipment is paid up front, and working capital
is locked up for ten years and returns at a heavily discounted value). Net of everything the
project earns more than its 14\% cost of capital.}

\refline{Berk, \S9.2 (Forecasting incremental earnings), pp.~250--254; \S9.3, pp.~254--260; \S9.4
(Sunk costs; fixed overhead expenses; the sunk cost fallacy), pp.~261--262.}

% =============================================================================
\section{Exercise 3 --- Replacing a one-year-old machine}

\noindent\textbf{Data (dollars).} Current machine: bought one year ago for $110{,}000$,
straight line over 11 years ($10{,}000$ per year), no salvage, EBITDA $20{,}000$ per year,
market value today $50{,}000$. New machine: price $150{,}000$, straight line over 10 years, no
salvage, EBITDA $40{,}000$ per year for 10 years. Other expenses identical. $\tau_c=0.45$,
$r=0.10$.

\begin{notation}
  \item $P_{new}=150{,}000$: price of the new machine; $P_{old}=110{,}000$: historical cost of the old one.
  \item $BV_{old}$: book value of the old machine today; $MV_{old}=50{,}000$: its market value today.
  \item $E_{new}=40{,}000$, $E_{old}=20{,}000$: EBITDA per year of each machine.
  \item $D_{new}$, $D_{old}=10{,}000$: annual depreciation of each machine.
  \item $\tau_c=0.45$; $r=0.10$; $\Delta$: difference ``replace'' minus ``keep''.
\end{notation}

\noindent\textbf{Assumptions.} The horizon is the next 10 years. The old machine has exactly 10
years of depreciable life left (11 minus 1) and would run, if kept, for those 10 years with no
salvage value, the same horizon as the new one. The EBITDA figures are each machine's
contribution, so the incremental EBITDA of replacing is the difference.

\step{1}{The old machine's purchase price is sunk, but its book value matters} The $110{,}000$
paid last year will not come back under either choice. It still matters for taxes, because it
determines the book value against which the sale is taxed. After one year of depreciation:
\begin{align}
  BV_{old} &= P_{old} - 1\times D_{old} \notag\\
           &= (110{,}000) - (10{,}000) \notag\\
           &= 100{,}000 . \label{eq:3bv}
\end{align}

\step{2}{After-tax cash from selling the old machine} Selling at $MV_{old}=50{,}000$ below book
value is a capital loss, which reduces taxes:
\begin{align}
  \text{Tax on sale} &= \tau_c\,(MV_{old}-BV_{old}) \notag\\
                     &= (0.45)\bigl((50{,}000)-(100{,}000)\bigr) \notag\\
                     &= -22{,}500 , \label{eq:3tax}
\end{align}
a tax \emph{saving} of $22{,}500$. The after-tax proceeds are
\begin{align}
  \text{After-tax sale} &= MV_{old} - \text{Tax on sale} \notag\\
                        &= (50{,}000) - (-22{,}500) \notag\\
                        &= 72{,}500 . \label{eq:3sale}
\end{align}

\step{3}{Incremental cash flow at $t=0$} Buy the new machine, sell the old one:
\begin{align}
  \Delta FCF_0 &= -P_{new} + \text{After-tax sale} \notag\\
               &= -(150{,}000) + (72{,}500) \notag\\
               &= -77{,}500 . \label{eq:3fcf0}
\end{align}

\step{4}{Incremental depreciation} The new machine depreciates
\begin{equation}\label{eq:3dnew}
  D_{new} = \frac{150{,}000}{10} = 15{,}000 \text{ per year},
\end{equation}
so replacing raises depreciation by
\begin{equation}\label{eq:3dd}
  \Delta D = (15{,}000) - (10{,}000) = 5{,}000 \text{ per year (years 1--10)}.
\end{equation}
The incremental EBITDA is
\begin{equation}\label{eq:3de}
  \Delta E = (40{,}000) - (20{,}000) = 20{,}000 \text{ per year}.
\end{equation}

\step{5}{Incremental FCF in years 1--10} No further capex or working capital, so
$\Delta FCF = \Delta E(1-\tau_c) + \tau_c\,\Delta D$ (after-tax EBITDA plus the extra depreciation
tax shield). Substituting \eqref{eq:3de} and \eqref{eq:3dd}:
\begin{align}
  \Delta FCF_t &= (20{,}000)(1-0.45) + (0.45)(5{,}000) \notag\\
               &= 11{,}000 + 2{,}250 \notag\\
               &= 13{,}250, \qquad t = 1,\dots,10 . \label{eq:3fcf}
\end{align}

\step{6}{Annuity factor at 10\%}
\begin{align}
  (1.10)^{10} &= 2.593742, \notag\\
  A(0.10,10) &= \frac{1-(2.593742)^{-1}}{0.10} = \frac{1-0.385543}{0.10} = 6.144567 . \label{eq:3af}
\end{align}

\step{7}{Incremental NPV} Substituting \eqref{eq:3fcf0}, \eqref{eq:3fcf} and \eqref{eq:3af}:
\begin{align}
  \Delta\mathrm{NPV} &= \Delta FCF_0 + \Delta FCF\cdot A(0.10,10) \notag\\
                     &= (-77{,}500) + (13{,}250)(6.144567) \notag\\
                     &= -77{,}500 + 81{,}415.51 \notag\\
                     &= 3{,}915.51 . \label{eq:3npv}
\end{align}

\answer{Yes, replacing is profitable: the incremental NPV is about \$3{,}915.51 $>0$ (incremental
IRR $\approx 11.2\% > 10\%$).}

\step{8}{Verification: value each alternative separately} \emph{Keep}: FCF of the old machine
$= 20{,}000(0.55) + 0.45(10{,}000) = 11{,}000 + 4{,}500 = 15{,}500$ per year, worth
$15{,}500\times6.144567 = 95{,}240.79$. \emph{Replace}: FCF of the new machine
$= 40{,}000(0.55) + 0.45(15{,}000) = 22{,}000 + 6{,}750 = 28{,}750$ per year, worth
$28{,}750\times6.144567 = 176{,}656.30$, minus the net outlay $77{,}500$, i.e.\ $99{,}156.30$. The
difference $99{,}156.30 - 95{,}240.79 = 3{,}915.51$ matches \eqref{eq:3npv}.

\remark{The tax saving on the sale is decisive. If one ignored it and used only the $50{,}000$
market price, the outlay would be $100{,}000$ and the NPV
$-100{,}000 + 81{,}415.51 = -18{,}584.49 < 0$, the opposite decision. Selling an asset below book
value creates a deductible loss worth $\tau_c(BV-MV)$; this is why book value, though built on a
sunk cost, still enters a cash-flow analysis.}

\intuition{The new machine doubles EBITDA and carries a larger depreciation tax shield; the cost is
a net 77,500 today once the old machine's resale value and the tax loss it generates are counted.
The extra after-tax cash, 13,250 a year for ten years, earns about 11.2\%, slightly above the 10\%
opportunity cost, so the replacement adds a modest 3.9 thousand of value.}

\refline{Berk, \S9.4 (Sunk costs; adjusting free cash flow: liquidation or salvage value and the
tax on capital gains; replacement decisions), pp.~261--265.}

\end{document}
